\documentclass[11pt,a4paper]{article}
\usepackage[T1]{fontenc}
\usepackage{lmodern,amsmath,amssymb,booktabs,longtable,array,geometry,hyperref}
\geometry{margin=23mm}
\hypersetup{hidelinks}
\emergencystretch=3em
\newcommand{\A}{\textbf{(a)}\ }
\newcommand{\B}{\textbf{(b)}\ }
\newcommand{\C}{\textbf{(c)}\ }
\newcommand{\D}{\textbf{(d)}\ }
\newcommand{\dd}{\mathrm d}
\newcommand{\file}[1]{\nolinkurl{#1}}
\newcommand{\hk}{\file{shiftcharge_bridge_hikami.py}}
\newcommand{\df}{\file{shiftcharge_bridge_deformation.py}}
\newcommand{\scs}{\file{shiftcharge_bridge_sector.py}}
\newcommand{\hy}{\file{shiftcharge_bridge_hierarchy.py}}
\newcolumntype{L}[1]{>{\raggedright\arraybackslash}p{#1}}
\title{Toward exact canonical good-Boussinesq scattering\\
Poisson contraction, soliton sectors and current deformations}
\author{Quantum Boussinesq solitons project\\PI Gustav Delius}
\date{4 November 2026}
\begin{document}
\maketitle
\begin{abstract}
\C A coordinated investigation tests the proposed quantum-model route beyond
the equation-level dictionary. The standard $W_3$ construction quantises the
second Gelfand--Dickey bracket, whereas the specified potential theory uses
the first. A singular large-spin-background contraction recovers the canonical
bracket with an explicit $\hbar$ normalization. Its multiplicative Baxter
arguments acquire finite additive shifts whose spacing in the soliton variable
$\ell$ is $9\hbar/2$. The inherited highest-weight annihilation conditions
are incompatible with a positive norm and the required real target fields.
The normalization and existence of the limiting Baxter operators remain open.

\C Independently, an anchored current deformation gives an exact classical
Hamiltonian family. It preserves isolated-soliton energy and momentum and
has first-order two-body action response $-g(P_1M_2-P_2M_1)$.
Its leading local density
identifies a dynamical matching parameter that one-body energy and vacuum
elasticity miss. A Gaussian infrared calculation shows that changing soliton
charges requires a nontrivial state dictionary even with arbitrarily cheap
finite-volume compensation.

\D The exact physical S-matrix and the seed phase coefficient remain
undetermined. The constructive outcome is an explicit candidate algebraic
contraction, a relative scattering deformation, and quantitative tests for the
physical representation and soliton states needed to complete the construction.
\end{abstract}

\section{Scope and research verdict}
Claims use the project's meanings of (a), literature result, (b), literature
heuristic, (c), project derivation or computation, and (d), unestablished
identification or research proposal. Conditional algebraic and Gaussian
results retain their stated scope. Every new displayed (c) formula is checked
by the script named in its subsection. The four scripts live in
\file{phase2/scripts/} and write only their adjacent JSON caches.

\C The work followed the proposed canonical real-form and one-soliton
dictionary test, an independent conserved-current deformation calculation,
and an independent sector audit. The two latter investigations cross-reviewed
the contraction, normalization and domain qualifications. A further
first-bracket conservation-law diagnostic was derived separately.
Existing packages, controlling audits and the previous
\file{shiftcharge_oneloop.tex} were used as context. The optional
\file{audit_good_renorm.tex} and \file{audit_defect_bootstrap.tex} were absent
from the repository inventory. No old script was run or imported.

\C The fixed target is
\begin{equation}
 H_0=\int\left[\frac{\pi^2+u^2+u_x^2}{2}-u^3\right]\dd x,\qquad
 u=\phi_x,\quad \{\phi(x),\pi(y)\}=\delta(x-y).
 \label{target}
\end{equation}
It gives $u_t=\pi_x$ and
$\pi_t=(u-3u^2-u_{xx})_x$. Its charges are
$M=\int u\,\dd x$, $\Pi=\int\pi\,\dd x$, and
$P=-\int\pi u\,\dd x$. The finite quadratic subtraction
$\lambda\hbar\int u^2\dd x$ remains a parameter.
\D Transferring source normal ordering to a physical target representation
is required before any value of $\lambda$ can be inferred.

\begin{center}
\small
\begin{tabular}{@{}L{4.0cm}L{5.7cm}L{5.3cm}@{}}
\toprule
Question & Result established & Remaining input\\
\midrule
Direct $W_3$ identification & \C Brackets disagree at the target vacuum &
\D A change of representation or a singular limit\\
Canonical contraction & \C Correct first bracket, energy and $\hbar$ scale &
\D Positive representation and operator convergence\\
Baxter equation & \C Finite argument shifts and the $9\hbar/2$ unit &
\D Vacuum-normalized functions and soliton eigenstates\\
$P\wedge M$ freedom & \C Explicit classical current-deformation direction &
\D Absolute seed scattering and quantum generator\\
Soliton state space & \C Infrared norm and finite-volume compensation tests &
\D Interacting dressing and scattering limits\\
Exact physical S-matrix & \D Undetermined &
\D The preceding physical inputs and channel closure\\
\bottomrule
\end{tabular}
\end{center}

\section{The canonical dictionary and its obstruction}
\subsection{Source bracket and independent reconstruction}
\A Bazhanov--Hibberd--Khoroshkin (BHK) explicitly identify their classical
algebra as the second Gelfand--Dickey structure. Their scalar Lax operator
is $D^3-UD+V-U_x/2-\lambda_{\rm sp}^3$.
Hikami supplies a lattice chiral algebra and a third-order Baxter relation.
The relevant full-text passages were read as recorded in
Section~\ref{sources}. \cite{BHK,Hikami}

\C (\hk) With coefficients placed to the left of $D=\partial_x$, the
second bracket is independently reconstructed as
\begin{align}
 (J_2)_{UU}&=2D^3-2UD-U_x,&
 (J_2)_{UV}&=-3VD-2V_x, \nonumber\\
 (J_2)_{VU}&=-3VD-V_x.&& \nonumber
\end{align}
\begin{equation}
 (J_2)_{VV}=-\frac{D^5}{6}+\frac{5U}{6}D^3+\frac{5U_x}{4}D^2
 +\left(\frac{3U_{xx}}4-\frac{2U^2}3\right)D
 +\frac{U_{xxx}}6-\frac{2UU_x}3 .
 \label{j2}
\end{equation}
The calculation starts from the source Miura fields
\begin{equation}
 U=a^2+b^2+2a_x,\qquad
 \sqrt3V=\frac23b^3-2a^2b-3ab_x-a_xb-b_{xx},
 \label{miura}
\end{equation}
and the current bracket $-\tfrac12D$ in each component. Composing their
Fr\'echet derivatives with this bracket and the formal adjoint reproduces
all four blocks of \eqref{j2}. Formal skew-adjointness is checked separately.

\C The exact real PDE dictionary is
\begin{equation}
 U=\frac14-\frac32u,\qquad V=\frac{\sqrt3}{4}\pi,\qquad
 \tau=\sqrt3\,t.
 \label{dictionary}
\end{equation}
The source flow $U_\tau=-2V_x$,
$V_\tau=U_{xxx}/6-2UU_x/3$ then gives \eqref{target}.
However the target canonical bracket in these variables is
\begin{equation}
 J_1=-\frac{3\sqrt3}{8}
 \begin{pmatrix}0&D\\D&0\end{pmatrix}.
 \label{j1}
\end{equation}
Its $UU$ bracket vanishes, whereas \eqref{j2} has a third derivative and
field-dependent terms. Scalar rescaling of time or $\hbar$ cannot identify
the displayed observable maps.

\subsection{How far the obstruction extends}
\C (\hk) At $U=1/4,V=0$,
\begin{equation}
 J_2=\operatorname{diag}\left(
 2D(D^2-\tfrac14),-\frac{D(D^2-1)(D^2-\tfrac14)}6\right),
 \qquad
 \frac{\det J_2}{\det J_1}
 =\frac{64}{81}(D^2-\tfrac14)^2(D^2-1).
 \label{det-ratio}
\end{equation}
The nonconstant determinant ratio excludes a finite-order differential
linear change with a finite-order differential inverse at this vacuum.
More generally, for constant backgrounds the Hermitian Fourier symbol is
\begin{equation}
 iJ_2(iq)=q
 \begin{pmatrix}
 2(q^2+U)&3V\\
 3V&(q^2+U)(q^2+4U)/6
 \end{pmatrix},
 \label{symbol}
\end{equation}
with determinant
\begin{equation}
 q^2\left[\frac{(q^2+U)^2(q^2+4U)}3-9V^2\right].
 \label{symbol-det}
\end{equation}
At the stated vacuum and $q>0$, $iJ_2(iq)$ is positive definite and
$iJ_1(iq)$ has one eigenvalue of each sign. Inertia under Hermitian congruence
excludes a real invertible translation-invariant linear Fourier map at this
same vacuum, including a nonlocal one.
\D Changes of vacuum, complex reality conditions and singular limits
require separate tests.

\section{A singular contraction with the correct quantum scale}
\subsection{Poisson pencil and Hamiltonian}
\C (\hk) A concrete escape from the fixed-vacuum obstruction is
\begin{equation}
 V_{\rm source}=V+\Lambda,\qquad
 \kappa_\Lambda=\frac{\sqrt3}{8\Lambda},\qquad
 \lim_{\Lambda\to+\infty}\kappa_\Lambda J_2(U,V+\Lambda)=J_1 .
 \label{contraction}
\end{equation}
The limit holds coefficient by coefficient for fields of finite size.
The shifted bracket has a large off-diagonal first-derivative term,
which survives this rescaling.

\C Let $F(U,V)=U^3+9V^2+3U_x^2/4$. The physical Hamiltonian has the exact
density representation
\begin{equation}
 H_0=\frac8{27}\int\left[F(U,V)-\frac3{16}U+\frac1{32}\right]\dd x.
 \label{energy-map}
\end{equation}
Direct substitution of \eqref{dictionary} verifies this identity, and
$J_1$ acting on its functional gradient gives physical $t$ evolution.
The required large-background subtraction obeys
\begin{equation}
 F(U,V+\Lambda)-18\Lambda(V+\Lambda)+9\Lambda^2=F(U,V).
 \label{subtraction}
\end{equation}
Thus a shifted higher integral with lower-integral subtractions supplies a
candidate Hamiltonian. The source's Hamiltonian time and its normalized
circle integrals must be converted explicitly.

\A BHK's spin-six quantum density contains
$:T^3:+9:W^2:+(c_{\rm V}-10):T_x^2:/32$.
\C With $T=(c_{\rm V}/24)U$ and
$W=i(-c_{\rm V}/24)^{3/2}(V+\Lambda)$, division by
$(c_{\rm V}/24)^3$ gives the derivative coefficient
\begin{equation}
 \frac34\left(1-\frac{10}{c_{\rm V}}\right).
 \label{quantum-density}
\end{equation}
\D The source regularized product is part of this algebraic prescription.
Its limit in a new physical representation, including the finite subtractions,
has yet to be constructed.

\subsection{Exact mode normalization and failure of the inherited vacuum}
\C (\hk) On a circle of circumference $2\pi$, use the source
$e^{-inx}$ convention and define
\begin{equation}
 U_n=\frac{24}{c_{\rm V}}\left(L_n-\frac{c_{\rm V}}{24}\delta_{n0}\right),
 \qquad
 V_n=\frac{W_n}{i(-c_{\rm V}/24)^{3/2}}-\Lambda\delta_{n0}.
 \label{modes}
\end{equation}
The exact mixed source commutator gives
\begin{equation}
 [U_n,V_m]=\frac{24}{c_{\rm V}}(2n-m)
 \left(V_{n+m}+\Lambda\delta_{n+m,0}\right).
 \label{mixed}
\end{equation}
Choosing
\begin{equation}
 c_{\rm V}(\Lambda)=-\frac{128\sqrt3\pi\Lambda}{\hbar}
 \label{central-scaling}
\end{equation}
makes its central term precisely
$-3\sqrt3\hbar n\,\delta_{n+m,0}/(16\pi)$, the Fourier form of
\eqref{j1}. The $UU$ and $VV$ commutators vanish in the coefficientwise
contraction at fixed mode indices and finite rescaled fields.
This establishes an algebraic $\hbar$ dictionary.

\C The required physical involution is
$U_n^\dagger=U_{-n}$, $V_n^\dagger=V_{-n}$. With the factor of $i$ in
\eqref{modes}, this requires $L_n^\dagger=L_{-n}$ and
$W_n^\dagger=-W_{-n}$ when $c_{\rm V}<0$.
Retaining the source highest-weight annihilation conditions
$U_n|0\rangle=V_n|0\rangle=0$ for $n>0$ then gives the limiting Gram matrix
\begin{equation}
 G_n\longrightarrow
 \begin{pmatrix}0&K_n\\K_n&0\end{pmatrix},
 \qquad K_n=-\frac{3\sqrt3\hbar n}{16\pi}.
 \label{gram}
\end{equation}
Its eigenvalues have opposite signs. Already the exact $UU$ entry is
$48n(n^2+U_0)/c_{\rm V}$, which is negative at $n=1,U_0=1/4$.
At $\Lambda=1000,\hbar=1$ the two independently evaluated level-one
eigenvalues are approximately $-0.1034244254$ and $0.1033239228$.
This excludes the inherited annihilation conditions with the required
target reality structure as a positive vacuum.

\C For constant real Miura currents, \eqref{miura} also implies
\begin{equation}
 4U^3-27V^2=4a^2(a^2-3b^2)^2\geq0.
 \label{discriminant}
\end{equation}
At $U=1/4$ it requires $|V|\leq1/(12\sqrt3)$.
The large-$\Lambda$ path therefore leaves the real constant-Miura sector.
\D A different polarization or a different representation could still
realise the contracted canonical algebra. The existing CFT eigenstates
do not yet supply that construction.

\subsection{Baxter arguments become additive}
\A BHK identify the transfer-matrix spectral parameter as
$t_{\rm sp}=\lambda_{\rm sp}^3$, and give the exact relation
\begin{equation}
 Q(t_{\rm sp}q^3)-T(t_{\rm sp}q^{1/2})Q(t_{\rm sp}q)
 +\overline T(t_{\rm sp}q^{-1/2})Q(t_{\rm sp}q^{-1})
 -Q(t_{\rm sp}q^{-3})=0.
 \label{source-baxter}
\end{equation}
Here $q=e^{i\pi\mathfrak g}$ and
$c_{\rm V}=50-24(\mathfrak g+\mathfrak g^{-1})$.
The bar labels the second transfer operator.

\C (\hk) Select the small positive $\mathfrak g_\Lambda$ root and
put $t_{\rm sp}=\Lambda+z$. Equations \eqref{central-scaling} and
\eqref{source-baxter} give
\begin{equation}
 \eta=\frac{\sqrt3\hbar}{16},\qquad
 A_r(z)=q_\Lambda^r(\Lambda+z)-\Lambda
 \longrightarrow z+ir\eta .
 \label{arguments}
\end{equation}
For $Q_\Lambda(z)=Q(\Lambda+z)$, and analogously for the transfer
operators, the finite-$\Lambda$ relation is exactly
\begin{align}
 Q_\Lambda(A_3(z))
 &-T_\Lambda(A_{1/2}(z))Q_\Lambda(A_1(z))\nonumber\\
 &+\overline T_\Lambda(A_{-1/2}(z))Q_\Lambda(A_{-1}(z))
 -Q_\Lambda(A_{-3}(z))=0 .
 \label{finite-baxter}
\end{align}
The argument limits are checked symbolically and with 80-digit arithmetic.
At $\Lambda=10^{20},\hbar=1,z=0.70$ the largest of the six sampled
argument errors is below $2.2\times10^{-21}$.

\D If compatible normalizations and domains give limits of the operators
or eigenfunctions, their prospective equation is
\begin{equation}
 Q(z+3i\eta)-T(z+i\eta/2)Q(z+i\eta)
 +\overline T(z-i\eta/2)Q(z-i\eta)-Q(z-3i\eta)=0.
 \label{limit-baxter}
\end{equation}
A $z$-dependent scalar renormalization of $Q_\Lambda$ changes the
coefficients. Establishing these normalized limits is a dynamical step.

\C The soliton roots
\begin{equation}
 \mu_\pm=\frac{-v/\sqrt3\pm k}{2},\qquad
 k^2+v^2=1,\qquad \ell=v(4k^2-1)
\end{equation}
satisfy
\begin{equation}
 \mu_+^3-\frac{\mu_+}4
 =\mu_-^3-\frac{\mu_-}4
 =z=-\frac{\ell}{12\sqrt3}.
 \label{spectral}
\end{equation}
Successive $Q$ arguments in \eqref{limit-baxter} are separated by
$2i\eta$, giving
\begin{equation}
 \Delta\ell=-12\sqrt3(2i\eta)=-\frac92i\hbar .
 \label{spacing}
\end{equation}
This agrees with the unit in the project's semiclassical string condition.
\D Identification of actual roots, physical poles or fusion amplitudes
requires the limiting eigenfunctions and their state interpretation.

\section{A dynamical realization of the antisymmetric direction}
\subsection{Currents and the anchored generator}
\A Generalized current deformations produce relative scattering dressings
under specified locality and scattering assumptions. The source treatments
retain finite-volume boundary terms and allow momentum-dependent charge
eigenvalues. Their full-text derivations motivate the following independent
Boussinesq calculation. \cite{DDY,CD}

\C (\df) The target conservation laws have densities and currents
\begin{equation}
 q_M=u,\quad j_M=-\pi,\qquad
 q_P=p=-\pi u,\quad
 j_P=\frac{\pi^2+u^2+u_x^2}{2}-2u^3-u u_{xx}.
 \label{currents}
\end{equation}
Both continuity equations are checked directly from \eqref{target}.
For a fixed left endpoint $\phi_-$, let $\chi=\phi-\phi_-$ and set
\begin{equation}
 X=-\int_{x<y}u(x)p(y)\,\dd x\dd y
   =\int\chi\pi u\,\dd x .
 \label{generator}
\end{equation}
The integrated current cross product is
\begin{align}
 {\cal D}&=q_Mj_P-j_Mq_P
 =-\frac12\pi^2u+\frac12u^3-2u^4-u^2u_{xx}+\frac12u u_x^2,\nonumber\\
 \int{\cal D}\,\dd x
 &=\int{\cal D}_{\rm bulk}\,\dd x-[u^2u_x],\nonumber\\
 {\cal D}_{\rm bulk}
 &=-\frac12\pi^2u+\frac12u^3-2u^4+\frac52u u_x^2 .
 \label{deformation-density}
\end{align}
For localized derivatives and vanishing currents the boundary term vanishes,
and $\partial_gH_g=\{H_g,X\}$ at $g=0$ gives this perturbation.
An independently checked local identity is
\begin{equation}
 \partial_t(\chi\pi u)=\partial_x(\chi j_P)-{\cal D}.
 \label{local-generator}
\end{equation}

\subsection{The exact classical family}
\C (\df) The anchored canonical flow is
\begin{equation}
 \partial_g\chi=\chi\chi_x,\qquad
 \partial_g\pi=\chi\pi_x.
\end{equation}
Its characteristic solution is
\begin{equation}
 x=y-g\chi_0(y),\quad
 \chi_g(x)=\chi_0(y),\quad
 \pi_g(x)=\pi_0(y),\quad
 u_g(x)=\frac{u_0(y)}{1-gu_0(y)} .
 \label{reciprocal}
\end{equation}
The chart requires $1-gu_0(y)>0$.
Writing $J=1-gu_0$, the variations at fixed $x$ obey
$\delta\chi_g=\delta\chi_0/J$ and
$\delta\pi_g=\delta\pi_0+g\pi'_0\delta\chi_0/J$.
Since $\dd x=J\dd y$, their wedge product preserves the canonical
symplectic form.

\C Transforming the Hamiltonian gives the all-order classical expression
\begin{equation}
 \boxed{
 H_g=\int\left[
 \frac{\pi^2(1-gu)}2+\frac{u^2}{2(1-gu)}
 -\frac{u^3}{(1-gu)^2}+\frac{u_x^2}{2(1-gu)^5}
 \right]\dd x .}
 \label{hg}
\end{equation}
Differentiation at $g=0$ reproduces \eqref{deformation-density}.
\D Quantum ordering, contact subtractions and a global self-adjoint
generator have yet to be constructed. The possible singularity of
\eqref{reciprocal} is a concrete limitation of the classical chart.

\C Within that anchored chart the field integrals obey
\begin{equation}
 T_gM=M,\qquad T_gP=P,\qquad T_g\Pi=\Pi+gP.
 \label{chargeflow}
\end{equation}
Here the left endpoint is a fixed c-number.
\D If it is promoted to a dynamical global shift variable, the complete
boundary symplectic problem must be solved. A uniform shift of $\phi$ and
$\phi_-$ together leaves the bilocal expression invariant. The formal
constant-shift vector field is also outside the fixed-endpoint tangent
space. Consequently \eqref{chargeflow} is used as a verified anchored
field-integral response, without asserting a global quantum commutator.

\subsection{One-body invariance and the centroid identity}
\C (\df) A solution of $H_g$ is obtained by applying $T_{-g}$ to a
solution of $H_0$. For an isolated soliton,
\begin{equation}
 x=y+g\chi_0(y,t),\qquad
 u_{\rm phys}=\frac{u_0}{1+gu_0},\qquad \pi_{\rm phys}=-v u_0 .
\end{equation}
Direct integration, separately from the canonical argument, gives
\begin{equation}
 M_g=M_0,\quad P_g=P_0,\quad E_g=E_0,\quad
 \Pi_g=\Pi_0-gP_0,\quad B_g=B_0+\frac{gM_0^2}{2},
 \qquad B=\int xu\,\dd x.
 \label{onebody-deform}
\end{equation}
The centroid retains its original velocity. The deformed Hamiltonian has
$\dot\phi=(1-gu)\pi$, so its integrated continuity identity is
\begin{equation}
 -M v_{\rm cm}=\int(1-gu)\pi\,\dd x=\Pi+gP.
 \label{deformed-ward}
\end{equation}
The original identity uses the original kinetic Hamiltonian.
This calculation explains how a proportional shift of $\Pi$ can accompany
an unchanged dispersion in a different Hamiltonian family.

\C On a classical one-soliton profile,
$\pi=-vu$, $v^2=1-k^2$, and $u_x^2=k^2u^2-2u^3$ imply
\begin{equation}
 {\cal D}_{\rm bulk}=3k^2u^3-7u^4,\qquad
 \int u^3\dd x=\frac{4k^5}{15},\quad
 \int u^4\dd x=\frac{4k^7}{35},\quad
 \int{\cal D}_{\rm bulk}\dd x=0 .
 \label{onebody-zero}
\end{equation}
The vacuum potential in \eqref{hg} has $V_g''(0)=1$ exactly.
For $g=O(\hbar)$, this direction has no quadratic vacuum term or
one-soliton energy shift at order $\hbar$. Its independent cubic kinetic
operator $-\pi^2u/2$ also distinguishes it from $\lambda u^2$.
Thus these one-body and vacuum measurements do not determine this
relative scattering deformation at that order.

\subsection{Two-body phase and an independent collision test}
\C (\df) Fix incoming ordering $x_1<x_2$ and outgoing ordering
$x_1>x_2$. For isolated symmetric solitons,
$X_j=-M_jP_j/2$. The separated generator has limits
\begin{equation}
 X_{\rm in}=X_1+X_2-M_1P_2,\qquad
 X_{\rm out}=X_1+X_2-M_2P_1 .
\end{equation}
Consequently
\begin{equation}
 X_{\rm out}-X_{\rm in}
 =-(P_1M_2-P_2M_1)=-Q_{12},\qquad
 \Delta I=-g\int{\cal D}\,\dd x\dd t=-gQ_{12}.
 \label{action-shift}
\end{equation}
The last equality is the first-order fixed-momentum action response.
It gives a leading semiclassical relative factor
\begin{equation}
 \exp\left[-\frac{ig}{\hbar}Q_{12}\right].
 \label{relative-factor}
\end{equation}
Taking $g=-\hbar c_{\rm CDD}$ realises the displayed order-one direction.
This is the full relative phase response. Computing only the response of
the old topological cross term would omit other dynamical and frame terms.

\C The checker constructs an independent exact two-soliton tau function
at $k_1=0.25,k_2=0.60$, both right-moving. It first tests the canonical field
equations, then integrates $X$ at $t=\pm2500$. It finds
\begin{align}
 X_{\rm in}&=-0.1292414735326871203679552509113167,\nonumber\\
 X_{\rm out}&=-0.08374454648958529813414045528563667,\nonumber\\
 X_{\rm out}-X_{\rm in}
 &=0.04549692704310182223381479562568000=-Q_{12}.
\end{align}
The incoming and outgoing factorizations each agree within $2\times10^{-47}$
against a declared $10^{-37}$ tolerance. A separate calculation at the
collision gives
\begin{equation}
 \int{\cal D}_{\rm bulk}(x,0)\,\dd x
 =-0.0004539822987511666070508801904128653,
\end{equation}
and verifies it against $-\dot X$ to better than $10^{-65}$.
This is a current-insertion check through a collision. The projected
functional determinant remains uncomputed.

\D If a self-adjoint renormalized $X$, a common domain and asymptotic
soliton states exist, the candidate unitary is $U_g=e^{igX/\hbar}$.
With $H_g=U_gH_0U_g^\dagger$ and a fixed asymptotic coordinate
identification, the relative quantum relation would be
\begin{equation}
 S_g(P_1,P_2)=
 e^{-igQ_{12}/\hbar}S_0(P_1,P_2).
 \label{conditional-quantum}
\end{equation}
Transforming the observables and asymptotic identification along with the
Hamiltonian gives the corresponding equivalent description.
Holding the laboratory coordinate, field and circumference fixed makes
\eqref{hg} a different Hamiltonian representative.

\D Equation \eqref{conditional-quantum} leaves the seed $S_0$ undetermined.
Specifying $g=0$ selects the original member of this family. It does not
calculate the coefficient called $c$ in an incomplete seed-phase ansatz.
The full functional continuum ambiguity also remains unclassified.

\subsection{Finite circumference retains a boundary contribution}
\C (\df) On an interval $[a,b]$, the anchored generator is
\[
 X=-\int_a^b p(y)\int_a^y u(x)\dd x\dd y .
\]
Using the primitive anchored at $a$ at each time, the continuity equations give
\begin{equation}
 -\dot X=\int_a^b{\cal D}\dd x+j_M(a)P-Mj_P(b).
 \label{interval}
\end{equation}
At periodic endpoints this reads
$\int{\cal D}=-\dot X+Mj_P(a)-j_M(a)P$.
A finite-interval numerical profile with nonzero endpoint currents checks
the complete identity independently to better than $10^{-65}$.
The reciprocal coordinate map sends a period $L$ with potential twist $M$
to $L-gM$.
\D Identifying $-\dot X$ with $\{H,X\}$ requires a boundary phase space in
which both functionals have Hamiltonian vector fields.

\C The corresponding algebraic ring-phase contribution is
\begin{align}
 \sum_l\Delta\delta_{jl}
 &=-\frac g\hbar(P_jM_{\rm tot}-M_jP_{\rm tot}),\nonumber\\
 \frac{P_jL}{\hbar}+\sum_l\Delta\delta_{jl}
 &=\frac{P_j(L-gM_{\rm tot})}{\hbar}
   +\frac{gM_jP_{\rm tot}}{\hbar}.
 \label{ring-shift}
\end{align}
\D Deriving a physical finite-volume spectrum requires a fixed
circumference and independently defined seam. The length response and the
charge-dependent twist must be kept in that dictionary.

\section{The soliton sector has a measurable infrared requirement}
\subsection{Conserved charges and the finite-ring zero mode}
\C (\scs) In the derivative-field algebra the first Hamiltonian
operator is $\left(\begin{smallmatrix}0&D\\D&0\end{smallmatrix}\right)$.
Constant gradients show that $M$ and $\Pi$ are central for compactly
supported derivative-field observables. In the potential algebra,
$\{\phi(x),\Pi\}=1$. With decaying bulk fields,
\begin{equation}
 \dot M=[\pi],\qquad
 \dot\Pi=[u-3u^2-u_{xx}].
\end{equation}
Endpoint variables need their own phase-space treatment.

\C The traveling branch has
\begin{equation}
 M=2k,\qquad \Pi=-2kv,\qquad
 P=\frac23k^3v=\frac{vM^3}{12},\qquad
 v=\epsilon\sqrt{1-M^2/4}.
 \label{fixed-sector}
\end{equation}
Fixing its kink height fixes the amplitude, leaving discrete velocity
signs. This bare classical family supplies no open momentum interval
inside one fixed-height sector.
\D Dressed field states can carry other momenta in that sector.
Their relation to derivatives along the classical $E(P)$ branch has to be
specified.

\C At fixed twist $M$ on a ring,
\begin{align}
 \phi(x)&=\Phi+\frac{Mx}{L}+\widetilde\phi(x),&
 \pi(x)&=\frac{\Pi}{L}+\widetilde\pi(x),\nonumber\\
 \Omega&=\dd\Pi\wedge\dd\Phi+
 \int\dd\widetilde\pi\wedge\dd\widetilde\phi .
 \label{ring-sector}
\end{align}
The tilded fields are periodic with zero mean.
The uniform configuration has
\begin{equation}
 H_{\rm unif}=\frac{\Pi^2+M^2}{2L}-\frac{M^3}{L^2},
 \qquad P_{\rm unif}=-\frac{M\Pi}{L}.
\end{equation}
\D Compactification of $\Phi$, allowed values of $M$, and quantum
boundary conditions remain additional data.

\subsection{Positive canonical oscillator polarization}
\C (\scs) The positive quadratic-vacuum representation can be written
explicitly. For nonzero ring modes with $q_n=2\pi n/L$, take
$\phi_n^\dagger=\phi_{-n}$,
$\pi_n^\dagger=\pi_{-n}$ and
$[\phi_n,\pi_m]=i\hbar\delta_{n+m,0}/L$. Define
\begin{equation}
 b_n=\sqrt{\frac{L}{2\hbar}}
 \left(\sqrt{\omega_n}\phi_n+\frac{i\pi_n}{\sqrt{\omega_n}}\right),
 \qquad \omega_n=|q_n|\sqrt{1+q_n^2}.
 \label{polarization}
\end{equation}
The independent oscillator checks give
$[b_n,b_m^\dagger]=\delta_{nm}$, $[b_n,b_m]=0$, and the inverses
\begin{equation}
 \phi_n=\sqrt{\frac{\hbar}{2L\omega_n}}(b_n+b_{-n}^\dagger),
 \qquad
 \pi_n=-i\sqrt{\frac{\hbar\omega_n}{2L}}(b_n-b_{-n}^\dagger).
\end{equation}
At a finite symmetric mode cutoff the quadratic Hamiltonian is
\begin{equation}
 H_2=\sum_{n\ne0}\hbar\omega_n
 \left(b_n^\dagger b_n+\frac12\right).
 \label{positive-energy}
\end{equation}
The modes $b_n$ and $b_{-n}$ are independent annihilation operators.
The source conditions that both $U_n$ and $V_n$ annihilate for $n>0$
would force both $b_n$ and $b_{-n}^\dagger$ to annihilate the state.
The latter condition is incompatible with a positive Fock vacuum.

\D The free canonical polarization is therefore available. The open
representation problem is to realise the limiting commuting densities and
transfer/Baxter operators with compatible ordering and domains in this
polarization, and to extend it to the charged soliton sectors.

\subsection{Exact Gaussian overlap test}
\C (\scs) Completing the square in the oscillator wavefunctions
gives overlap magnitude $e^{-N_{\rm diff}/2}$, where in the good free
vacuum
\begin{equation}
 N_{\rm diff}=\frac1{2\hbar}\int\frac{\dd q}{2\pi}
 \left[\omega(q)|\Delta\widehat\phi(q)|^2+
 \frac{|\Delta\widehat\pi(q)|^2}{\omega(q)}\right],
 \qquad \omega(q)=|q|\sqrt{1+q^2}.
 \label{occupation}
\end{equation}
For real initial data the cross term cancels between $q$ and $-q$.
At $q\ne0$, $\Delta\widehat\phi=\Delta\widehat u/(iq)$.
The global additive constant is a separate zero mode.
For the exponentially localized soliton profiles considered here,
$\Delta\widehat u(0)=\Delta M$ and
$\Delta\widehat\pi(0)=\Delta\Pi$. Hence
\begin{equation}
 \boxed{
 N_{\rm diff}(\mu)=
 \frac{(\Delta M)^2+(\Delta\Pi)^2}{2\pi\hbar}
 \log\frac1\mu+O(1).}
 \label{ir-log}
\end{equation}
The regulated overlap magnitude has power
$[(\Delta M)^2+(\Delta\Pi)^2]/(4\pi\hbar)$.
Changing these charges is therefore an infrared-divergent coherent
displacement of the quadratic-vacuum representation on the line.
\D This calculation does not determine an interacting infrared exponent
or establish the existence or failure of a dressed scattering operator.

\C The independent Fourier quadrature uses
\begin{equation}
 u_k(x)=\frac{k^2}{2}\operatorname{sech}^2(kx/2),\qquad
 \widehat u_k(q)=\frac{2\pi q}{\sinh(\pi q/k)},\qquad
 \widehat\pi_k(q)=-v\widehat u_k(q).
 \label{fourier}
\end{equation}
For right movers $k_1=0.31,k_2=0.53$ at $\hbar=1$, the predicted
occupation logarithm coefficient is
\begin{equation}
 A_N=0.04605009066875933660934658357463906 .
\end{equation}
Independent cutoff integrals at $\mu=10^{-6},10^{-12},10^{-24}$ converge
after subtracting $A_N\log(1/\mu)$ and reproduce its coefficient to
better than $2.0\times10^{-26}$. A fixed-amplitude translation by $1.7$
at $k=0.31$ instead has the finite limiting occupation
$0.005859161839779369150053652746685633$.

\subsection{Compensation in one finite-volume space}
\C (\scs) Periodize each localized profile and subtract its own
uniform $M/L$ and $\Pi/L$ pedestals. Both profiles now have the same
total charges $M=\Pi=0$. Their nonzero Fourier coefficients retain
\eqref{fourier}. For $q_n=2\pi n/L$,
\begin{equation}
 N_L=\frac1{\hbar L}\sum_{n\geq1}
 \left[
 \frac{\sqrt{1+q_n^2}}{q_n}|\Delta\widehat u(q_n)|^2+
 \frac{|\Delta\widehat\pi(q_n)|^2}{q_n\sqrt{1+q_n^2}}
 \right].
 \label{finite-overlap}
\end{equation}
These profiles define periodic initial data.
The exact nonlinear periodic traveling solution is a separate problem.
The independently summed doubling slopes approach $A_N$.

\begin{center}
\begin{tabular}{@{}rrr@{}}
\toprule
$L$ & $[N_{2L}-N_L]/\log2$ & Error from $A_N$\\
\midrule
256 & 0.0461099132406 & $5.98226\times10^{-5}$\\
512 & 0.0460650331524 & $1.49425\times10^{-5}$\\
1024 & 0.0460538254672 & $3.73480\times10^{-6}$\\
\bottomrule
\end{tabular}
\end{center}
\C No fitted logarithm coefficient enters this test. The occupation
quadratures and Fourier sums use $q\leq8$, where the exact hyperbolic-sine
factors suppress the tail. The many digits retained in JSON are arithmetic
output. The quoted convergence tests, rather than all those digits, are
the numerical accuracy claims. The logarithm survives charge compensation
in a common finite-volume oscillator space.

\C A separate construction places a smooth cloud far from the pulse,
\begin{equation}
 u_c(x)=\frac{\Delta M}{L}w((x-R)/L),\quad
 \pi_c(x)=\frac{\Delta\Pi}{L}w((x-R)/L),\quad
 \int w=1 .
\end{equation}
It has fixed compensating charges and
\begin{align}
 H_c&=\frac{(\Delta M)^2+(\Delta\Pi)^2}{2L}I_2
 -\frac{(\Delta M)^3}{L^2}I_3
 +\frac{(\Delta M)^2}{2L^3}I_d,\nonumber\\
 P_c&=-\frac{\Delta M\Delta\Pi}{L}I_2 ,
 \label{cloud}
\end{align}
where $I_2=\int w^2$, $I_3=\int w^3$, and $I_d=\int(w')^2$.
For $w=e^{-y^2}/\sqrt\pi$ the moments are
$I_2=I_d=1/\sqrt{2\pi}$ and $I_3=1/(\sqrt3\pi)$.
Direct spatial quadrature checks the complete charges and energy at two
widths. Thus vanishing compensation energy and a controlled state
normalization are separate requirements.

\D The quantum dictionary should identify operators intertwining
boundary-charge representations or a finite-volume family with controlled
soft compensation. The zero mode, nonzero phonon modes, and seam
normalization need explicit treatment before extracting a scattering phase.

\section{A first-bracket quantum-integrability diagnostic}
\C (\hy) An additional conserved density can be derived directly
in canonical fields, without using the $W_3$ representation. Let
\begin{equation}
 H_s=\int\left[\frac{\pi^2+s u^2+u_x^2}{2}-u^3\right]\dd x .
\end{equation}
Solving a weighted polynomial ansatz by setting both Euler derivatives of
its time derivative to zero gives, modulo a multiple of $H_s$,
\begin{align}
 {\cal Q}_7&=\int\left[
 \pi_x^2+u_{xx}^2-3u\pi^2+3u^4-9u u_x^2
 +s(u_x^2-u^3)\right]\dd x,\nonumber\\
 \{H_s,{\cal Q}_7\}&=0 .
 \label{q7}
\end{align}
This holds with appropriate vanishing boundary flux or periodic fields.
The script solves for the coefficients independently, then verifies
the final conservation law. Direct soliton integration at $s=1$ gives
\begin{equation}
 {\cal Q}_7[u_k,-vu_k]=\frac{4k^5(20k^2-21)}{105},
 \label{q7sol}
\end{equation}
with independent 70-digit quadrature at three amplitudes.

\C A finite-dimensional Weyl-symbol degree observation applies to
linear canonical regulators. $H_s$ is at most cubic in the coordinate
variables and quadratic in momenta, with no mixed kinetic terms.
The charge in \eqref{q7} is at most quadratic in momenta.
The only potentially nonzero third derivative of $H_s$ is purely in
coordinates, whereas the corresponding third momentum derivative of
$\mathcal Q_7$ vanishes. All higher odd Moyal contractions vanish by degree.
Thus for these regulated polynomial symbols
\begin{equation}
 [\operatorname{Op}_{\rm W}(H_s),
   \operatorname{Op}_{\rm W}({\cal Q}_7)]
 =i\hbar\operatorname{Op}_{\rm W}
   \bigl(\{H_s,{\cal Q}_7\}_{\rm regulator}\bigr).
 \label{weyl}
\end{equation}

\C The regulator's classical bracket need not vanish.
For the naive periodic three-site central-difference discretization,
at $\phi=(1,2,4)$ and $\pi=(2,-1,3)$ with $s=1$, it is exactly
$-75/8$. This rejected regulator is retained as an explicit failure.
\D Constructing renormalized continuum operators still requires a
compatible regulator, domain and ordering prescription. The finite
Weyl identity supplies a diagnostic for that work. Equation \eqref{q7}
also persists for arbitrary $s$, so this conservation-law test alone
cannot select the finite elastic subtraction.

\section{What can now be attempted toward the exact amplitude}
\D The first proposed assignment has produced a precise obstruction and
a candidate replacement. The remaining programme is organized by observable
tests rather than by an assumed scalar-factor ansatz.

\begin{enumerate}
\item Realise the limiting commuting densities in the explicit positive
canonical polarization \eqref{polarization}, replacing the inherited
highest-weight annihilation conditions. Determine the finite normal-ordering
terms and common domains. Equations \eqref{j1}, \eqref{energy-map},
\eqref{central-scaling} and \eqref{q7} provide explicit checks.
\item Obtain vacuum-normalized limits of the transfer and Baxter operators
at $t_{\rm sp}=\Lambda+z$. Establish convergence or identify an obstruction.
The six finite shifts in \eqref{arguments} and the spectral relation
\eqref{spectral} fix the proposed canonical scale.
\item Identify one-soliton states or intertwiners with the correct $M,\Pi,P$
and classical profile. Test their infrared normalization against
\eqref{ir-log} and \eqref{finite-overlap}, keeping the interacting nature
of the final exponent as a separate question.
\item Derive two-excitation finite-volume equations in that same
representation. Separate the dispersion, soft-state normalization,
potential seam and circumference response using \eqref{interval} and
\eqref{ring-shift}. Then extract the phase and compare several kinematic
points with a boundary-compatible continuum calculation.
\item Match the coefficient of \eqref{deformation-density} in fixed
canonical fields and physical coordinates. It is the explicit local
direction associated with the relative $P\wedge M$ factor.
This matching supplements the finite quadratic subtraction.
\end{enumerate}

\C The resonant-channel audit already requires care with heavy classical
solitons. \D The light branch supplies the initial candidate sector.
Closure under quantum scattering, including phonon and number-changing
channels, must follow from the constructed states and conserved operators.

\D At this stage there is no extracted physical $S_0$, no physical
two-soliton determinant, and no values of $c_{\rm good}(\lambda)$ or
$c_{\rm bad}$. The new result is a relative deformation law and a canonical
functional-equation route with its representation obstruction identified.
The earlier Toda comparison remains method validation with its audited
dictionary qualifications.

\section{Verification coverage and failures recorded}
\small
\begin{longtable}{@{}L{4.0cm}L{6.4cm}L{4.6cm}@{}}
\toprule
Claim group & Independent verification & Script\\
\midrule\endhead
Second versus first bracket & Miura Fr\'echet reconstruction, formal adjoints,
Fourier inertia and discriminant & \hk\\
Canonical contraction & PDE, energy, bracket coefficients and exact mixed modes &
\hk\\
Representation and Baxter scale & Exact finite-parameter Gram matrices,
80-digit argument limits and spectral polynomial reduction & \hk\\
Current deformation & Continuity laws, reciprocal Jacobian, symplectic form,
exact classical Hamiltonian and one-body integrals & \df\\
Scattering and ring response & Exact two-soliton field equations, separated
and collision quadratures, finite-interval boundary currents & \df\\
Soliton sectors & Positive oscillator algebra, Gaussian integral, exact Fourier transform, logarithm
coefficient, independent finite-ring sums and cloud quadratures & \scs\\
Canonical extra charge & Independent coefficient solve, Euler derivatives,
soliton quadrature and rejected lattice regulator & \hy\\
\bottomrule
\end{longtable}
\normalsize

\C The symbolic checks use SymPy. Numerical checks use 70 or 80 decimal
digits. The four scripts together run in about one minute on the
task machine, with the two-soliton deformation quadrature taking most of
the time. The deformation checker records 49 symbolic and numerical
checks, including independent charge quadratures and exact-tau
collision tests. Each script ends by printing \texttt{ALL CHECKS PASSED}.
JSON files distinguish computed quantities from unconstructed physical
observables. Numerical agreement certifies the stated tests.

\begin{enumerate}
\item \C Direct same-vacuum identification fails the $UU$ bracket test.
The stronger real Fourier-map attempt fails the inertia test.
\item \C The large-background contraction leaves the real constant-Miura
region. The inherited highest-weight annihilation conditions fail positivity
under the required target involution.
\item \C The bare traveling-soliton family has no open momentum branch at
fixed kink height. Low-energy charge compensation retains an infrared
logarithm in its Gaussian state normalization.
\item \C One-body dispersion and vacuum elasticity miss the current
deformation at order $\hbar$. This records failure of those observables to
fix the relative two-body direction.
\item \C The naive three-site canonical regulator fails the extra-charge
commutation test, with exact bracket $-75/8$.
\item \C Promoting the anchored charge response to a global quantum
commutator requires additional boundary data. The report retains that
qualification explicitly.
\item \C The Belov--Chaltikian HTML retrieval failed. The local full PDF
supplied the required passages. BHK's printed equation (2.24) repeats
$I_4$ for the spin-six density, and the cubic printed after equation (2.29)
disagrees with its displayed Lax operator. The source pages were visually
checked. All spectral calculations here use equation (2.26) and direct
polynomial substitution.
\item \C No Python implementation assertion failed and no numerical
tolerance was weakened. Two document-authoring submissions initially
failed because LaTeX dollar-brace syntax was interpreted by the JavaScript
template wrapper. Rephrasing that math command fixed the wrapper before
the affected text was written. The first PDF build exposed two overfull
lines, repaired by splitting the operator display and the interval formula.
A draft numerical-sample label was corrected from $z=0.37$ to the actual
checker input $z=0.70$ during independent review. The computed bound was
unchanged. Review also restricted the logarithm's bounded-remainder
statement to the exponentially localized profiles actually studied and
stated the interval identity in terms of $-\dot X$ before any boundary
Poisson identification.
\end{enumerate}

\section{Sources and reading depth}\label{sources}
\A Relevant full-text passages were read for each external reference below.
The reading did not reproduce every proof in the cited papers.

\begin{itemize}
\item Hikami \cite{Hikami}, sections 1--2, the Baxter result and analytic
assumption in section 3, sections 4--5, conclusion and appendices A--B.
Its inversion identity labelled unitarity needs a specified adjoint before
being used as physical Hilbert-space unitarity.
\item BHK \cite{BHK}, introduction, sections 2.1--2.2 and 3.2--3.3,
the spectral convention in section 4.2 and equation (1.11).
Printed pages 8--10 were visually inspected. The current work independently
checks the target normalization and the required reality structure.
\item Belov--Chaltikian \cite{BC}, sections 1--5, including the lattice
brackets, classical bi-Hamiltonian construction and continuum discussion.
\item Doyon--Durnin--Yoshimura \cite{DDY}, sections 2--3, appendix B and
the boundary formulas in appendix C.
\item Cardy--Doyon \cite{CD}, sections 2.2, 3.2--3.3 and 3.6.
\end{itemize}

{\raggedright
\C Internal inputs include the README Status, the ranked programme in
\file{report.tex}, \file{good_renorm.tex},
\file{shiftcharge_oneloop.tex}, \file{audit_resonant_channels.tex},
\file{audit_charges_action.tex}, \file{audit_projected_measure.tex},
\file{audit_det_projected.tex}, and the controlling errata.
The new scripts and this report use the prefix \file{shiftcharge}.
Existing files and git state were left untouched.\par}

\begin{thebibliography}{9}
\raggedright
\bibitem{Hikami}
K.~Hikami, \emph{Baxter Equation for Quantum Discrete Boussinesq Equation},
Nucl.\ Phys.\ B \textbf{604} (2001) 580--602,
\href{https://arxiv.org/abs/nlin/0102021}{arXiv:nlin/0102021}.

\bibitem{BHK}
V.~V.~Bazhanov, A.~N.~Hibberd and S.~M.~Khoroshkin,
\emph{Integrable structure of $W_3$ conformal field theory, quantum
Boussinesq theory and boundary affine Toda theory},
Nucl.\ Phys.\ B \textbf{622} (2002) 475--547,
\href{https://arxiv.org/abs/hep-th/0105177}{arXiv:hep-th/0105177}.

\bibitem{BC}
A.~A.~Belov and K.~D.~Chaltikian,
\emph{Lattice analogues of $W$-algebras and Classical Integrable Equations},
\href{https://arxiv.org/abs/hep-th/9303166}{arXiv:hep-th/9303166}.

\bibitem{DDY}
B.~Doyon, J.~Durnin and T.~Yoshimura,
\emph{The Space of Integrable Systems from Generalised
$T\overline T$-Deformations},
SciPost Phys.\ \textbf{13} (2022) 072,
\href{https://arxiv.org/abs/2105.03326}{arXiv:2105.03326},
\href{https://doi.org/10.21468/SciPostPhys.13.3.072}{doi:10.21468/SciPostPhys.13.3.072}.

\bibitem{CD}
J.~Cardy and B.~Doyon,
\emph{$T\overline T$ deformations and the width of fundamental particles},
JHEP \textbf{04} (2022) 136,
\href{https://arxiv.org/abs/2010.15733}{arXiv:2010.15733},
\href{https://doi.org/10.1007/JHEP04(2022)136}{doi:10.1007/JHEP04(2022)136}.
\end{thebibliography}
\end{document}
